\clearpage
\fancyhead[L]{\small BODE. ASTRONOMISCHES JAHRBUCH.}
\SourceLeaf{509}
\begin{center}\small Göttingische gelehrte Anzeigen. No. 155. pp. 1545..1552. 26 September 1808.\\[3pt]\rule{.72\linewidth}{.3pt}\end{center}

Available from the author and through \textsc{Fr. Braunes} on commission. \emph{Astronomical Yearbook for the Year} 1810, \emph{with a Collection of the Latest Papers and Reports Relating to the Astronomical Sciences. By} \textsc{J. E. Bode}. 1807. 268 pages in octavo and one engraved copperplate. Berlin.

This esteemed yearbook is indispensable to every astronomer. Its first part, which contains the astronomical calendar, is arranged as it was in the preceding volume. The worthy editor, however, has always worked tirelessly to make it more useful and has paid attention to proposals and requests to that end. We therefore take the liberty of presenting some such requests, which we know for certain several other astronomers have long shared. The detailed explanation of the calendar and its use has not been reprinted since the volume for 1800. We find this very commendable. The fairly long article would take up space that could be put to better use, and \emph{most} of its contents can in any case be assumed familiar to anyone who wants to observe the sky. Some parts of the explanation concern arbitrary conventions and matters that are not self-evident. These really should become a regular feature or be incorporated in some other way into the calendar. For example, no one would guess that all the planetary positions are given for midnight in true Berlin time, except those of Venus and Mercury, whose positions are for noon. Anyone unable to consult the volume for 1800 remains in doubt about this. The same applies to the occultations of fixed stars by the Moon, which are given for true time rather than the mean time now used by almost all astronomers. With the apparent obliquity of the ecliptic, the nutation is always given with its sign reversed. This convention hardly explains itself. One would naturally expect to subtract the nutation from the apparent obliquity, which includes it, to obtain the mean obliquity. With the entry in the yearbook the reverse is true. Nor can we approve of Mr \textsc{Bode} still giving the obliquity of the ecliptic from Mr \textsc{von Zach}’s older solar tables. According to the newer observations that English, French and Italian astronomers have been making for many years, it is $8''$ too large. Consequently all the Sun’s declinations in the yearbook are also too large and, if relied on, give incorrect altitudes of the pole depending on the season. It would also be desirable for Mr \textsc{Bode} to replace the inconvenient practice of giving times in morning and afternoon hours with the convention astronomers use. We know of several instances in which stellar occultations and eclipses of Jupiter’s satellites were missed solely because the times were given in civil time. Following a request made by a reviewer in issue 41 of the Literaturzeitung for 1797, Mr \textsc{Bode} has included heliocentric longitudes and latitudes in his planetary positions since 1800 and has left out the right ascensions in exchange. We have never been satisfied with this. We believe that geocentric right ascension and declination are the most essential quantities in an astronomical ephemeris. We would sooner do without all the other planetary columns than these two indispensable quantities. Every astronomer would surely welcome the addition of the distances from Earth on several counts, since these distances are indispensable in determining aberration. The reasons that reviewer had for requesting the heliocentric positions do not seem all that signif\SourceLeaf{510}icant to us. We believe that little use can be made of these two columns, especially since Mr \textsc{Bode} gives the oppositions and conjunctions of the planets among the monthly phenomena, and these could if need be be derived just as easily from the geocentric positions. Many years of using astronomical ephemerides have prompted further requests that we could add here. They would, however, be difficult to carry out without more extensive changes to the form and arrangement of the ephemerides. We therefore restrict ourselves to one request. By granting it the distinguished editor would greatly oblige all those who do much work on planetary and cometary computations and for whom calculating the positions of the Sun always takes a great deal of time. In very many cases that expenditure of time could be avoided altogether if the distances of the Sun from the Earth, or their logarithms, were given for every day of the year, even if necessary with one decimal place fewer than Mr \textsc{Bode} supplied in the last four volumes. The editor could make up for the slight increase in work by leaving out the seconds in the Moon’s longitudes and latitudes. They take an enormous amount of work to calculate and, we believe we can say, are little used if at all.

Our remarks on the arrangement of the yearbook have been prompted solely by our sincere wish to see it made still more useful. We now turn to the papers that give the astronomical yearbook lasting value. 1) Investigation of the true elliptical motion of the comet of 1769, by \textsc{Fr. W. Bessel}. This paper was submitted in competition for a prize of 30 Friedrichsd’or that an anonymous person had offered for the best astronomical paper or the most remarkable discovery. Mr \textsc{Bessel} was awarded only one half, and the other went to Mr \textsc{Huth} as co-discoverer of the two comets of 1805. Several astronomers had already worked on the elliptical orbit of that comet, and Father \textsc{Asklepi} had investigated it most carefully. The subject certainly deserved to be taken up again. The extensive reforms of our star catalogues and solar tables especially invited a new investigation, since with their help one could hope for much more reliable results. Mr \textsc{Bessel} has now treated the subject so thoroughly and so well that there is nothing left to wish for. If we may grant that the fundamental positions which Mr \textsc{Bessel} derived from a new and extremely careful analysis of all the good observations are reliable to within $5''$ in right ascension and declination, the comet’s orbital period is bounded by $1691\frac{3}{5}$ and 2673 years. Mr \textsc{Bessel}’s procedure for calculating the effect of refraction deserves particular attention. It is a pity that several misprints have completely disfigured the relevant formulas, so they cannot be used at all without consulting a later detailed article on the subject in the Mon. Corresp. 2) On the geographical longitude of Havana, by Mr \textsc{Oltmanns}. From the observations of \textsc{Churrucca}, \textsc{Robredo}, \textsc{v. Humboldt} and \textsc{Galiano}, the mean is $5^{\mathrm u}\,38^{\mathrm m}\,51^{\mathrm s}5$ west of Paris. 3) Tables for reducing the altitudes of Polaris to the meridian altitude, for the latitude of Berlin. 4) Apparent variations in the brightness of Algol for the years 1808, 1809, 1810, by Professor \textsc{Wurm}. 5) Third and fourth lists of the comparative brightness of the fixed stars, by Mr \textsc{Herschel}. 6) On the true geographical longitude of the degree of latitude measured in Peru, by Mr \textsc{Oltmanns}. From observations by \textsc{Humboldt}, \textsc{Ulloa} and the French academicians, the mean longitude of Quito is $5^{\mathrm u}\,24^{\mathrm m}\,20^{\mathrm s}$. The latest volume of the Connoissance des tems for 1809 still gave $5^{\mathrm u}\,21^{\mathrm m}\,0^{\mathrm s}$, following \textsc{Bouguer}. 7) Tables for calculating the annual change in the right ascension and declination of the fixed stars. Several errors have crept into this article. The action of the Sun and the Moon does not produce the annu\SourceLeaf{511}al change of $50''15$ in longitude, further increased by $0''135$ by the change in the obliquity of the ecliptic, as the author puts it. The action of the Sun and Moon produces the so-called lunisolar precession. The shift of the ecliptic slightly reduces it. That shift has the same cause as the decrease in obliquity, but the two must not be confused. The shift brings the precession down to $50''15$ or, according to more precise determinations, to $50''10$, the observed value. The coefficients of the variable part of the precession in right ascension and the formula for the change in declination should be exactly alike. The rule for calculating the change in the annual variation is also wrong. 8) A contribution to methods for using a series of lunar distances to calculate geographical longitudes, by Mr \textsc{Jabbo Oltmanns} (the difference between the apparent and true distances is calculated for three instants and determined from them by interpolation for all the observations). 9) Astronomical observations in Prague in 1806, by Messrs \textsc{David} and \textsc{Bittner}. 10) A method of determining the geographical positions of terrestrial objects from observed azimuths, angles of elevation and relative heights, with some auxiliary tables for calculating barometric height measurements, by Mr \textsc{Jabbo Oltmanns}. The procedure admittedly offers little accuracy, but it need not be dismissed in entirely unknown regions where other observations are lacking. 11) On \textsc{Troughton}’s tubular pendulum, by Mr \textsc{Schnitter} of Aachen. The yearbook for 1808 described \textsc{Troughton}’s invention in such a way that even the inventor’s own figures for the expansion of brass and steel did not yield the correct compensation. Mr \textsc{Schnitter} deserves thanks for drawing attention to this and to the fact that perfect compensation requires all the rods that expand downwards to be made of steel. The reviewer, however, has \textsc{Troughton}’s original article (\textsc{Nicholson}’s philosophical journal, vol. IX) before him and believes he must point out that the error lies solely in the yearbook’s incorrect translation. \textsc{Troughton}’s instructions are exactly as they must be according to Mr \textsc{Schnitter}’s account. 12) Astronomical observations and remarks by Mr \textsc{van Beeck Calkoen}. 13) Observation of the occultation of $\alpha$ in Scorpius on 2c.\ March 1805 and of the solar eclipse of 16 June 1806 on the island of Leon near Cadiz, by Mr \textsc{Canelas}. 14) Astronomical observations at Vienna in 1806, by Mr \textsc{Triesnecker}. 15) Discovery and observation of a fourth new planet by Dr \textsc{Olbers}. 16) Observation and calculation of the orbit of the comet of 1806, observations of Vesta and Juno, and stellar occultations, by Mr F. W. \textsc{Bessel}. 17) Measurement of the apparent size of Vesta, by Mr \textsc{Schröter} (on 26 April 1807 the apparent diameter was $0''488$, which gives a true diameter of 68 miles). 18) Observation of Vesta and calculation of the elements of its orbit and of its motion for 1808, by Dr \textsc{Gauss}. 19) Observations of Vesta at Berlin, by the editor. 20) Astronomical observations and remarks by Mr \textsc{Fritsch} in Quedlinburg. The observations of the planets, including those of Juno in 1807 when it was very difficult to observe, do Mr \textsc{Fritsch} credit, although they are of no use to the astronomer for lack of adequate time determinations. 21) On a mercury pendulum, by Mr \textsc{Th. Blacker} in London. 22) Astronomical observations by Mr \textsc{Lalande}. 23) Observations on the climates and atmosphere of Saturn, by Dr \textsc{Herschel}. The changes in the colour of Saturn’s polar regions noted by Dr H. appear to be periodic, and Mr H. regards them as proof of the existence of an atmosphere. Astronomical observations at Berlin Observatory in 1806. The upheavals of war in the autumn of that year caused various disruptions and interruptions. 25) Motion of Pallas and Juno in 1808 (calculated in advance). 26) On observed differences in the apparent magnitudes of some stars, by Dr \textsc{Koch} in Danzig. 27) Astronomical observa\SourceLeaf{512}tions in Spain (drawn from several Spanish publications). 28) Some remarks on physical astronomy, by Mr \textsc{Huth} in Frankfurt an der Oder. 29) A contribution to the determination of geographical longitude, by Mr \textsc{Oltmanns}. 30) Count \textsc{von Platen}’s proposed method for solving an astronomical problem rests essentially on converting products of sines into sums, as astronomers used to do before logarithms were invented. It is now more convenient to calculate with logarithms in most cases. The reviewer finds the proposed method inconvenient in precisely the cases the author discusses. 31) A proposed method of determining horizontal refraction from geographical longitude, by Mr \textsc{Oltmanns} (using the Moon’s distances from the Sun low in the sky). 32) Mr \textsc{Fischer}’s supplement to the paper on the best shape for objective mirrors in the yearbook for 1806. Here Mr F. withdraws his mistaken claim that spherical mirrors are preferable to parabolic ones. 33) Various astronomical observations and reports. One report deserves attention. Mr \textsc{Hecker} observed a fifth-magnitude star at Remplin near $\alpha$ in Equuleus on 17 August 1804. It was never seen again. We only wonder why Mr \textsc{Hecker} did not report this immediately, when it might have been useful.
